% настройки polyglossia
\setdefaultlanguage{russian}
\setotherlanguage{english}

% локализация
\addto\captionsrussian{
	\renewcommand{\partname}{Глава}
	\renewcommand{\contentsname}{Содержание}
	\renewcommand{\figurename}{Рисунок}
}

% основной шрифт документа

% перечень использованных источников
\addbibresource{../report/references.bib}

% оформление презентации
\usetheme{metropolis}
\usecolortheme{seagull}
\beamertemplatenavigationsymbolsempty

% Кириллические шрифты: системные при наличии, иначе файлы GNU FreeFont из TeX Live.
\IfFontExistsTF{CMU Serif}
    {\setmainfont{CMU Serif}}
    {\IfFontExistsTF{Times New Roman}
        {\setmainfont{Times New Roman}}
        {\setmainfont{FreeSerif.otf}}}
\IfFontExistsTF{Arial}
    {\setsansfont{Arial}\newfontfamily\cyrillicfontsf{Arial}[Script=Cyrillic]}
    {\setsansfont{FreeSans.otf}\newfontfamily\cyrillicfontsf{FreeSans.otf}[Script=Cyrillic]}
\IfFontExistsTF{Courier New}
    {\setmonofont{Courier New}\newfontfamily\cyrillicfonttt{Courier New}[Script=Cyrillic]}
    {\setmonofont{FreeMono.otf}\newfontfamily\cyrillicfonttt{FreeMono.otf}[Script=Cyrillic]}

% настройка ссылок и метаданных документа
\hypersetup{unicode=true,colorlinks=true,linkcolor=red,citecolor=green,filecolor=magenta,urlcolor=cyan,
	pdftitle={\docname},
	pdfauthor={\studentname},
	pdfsubject={\docname},
	pdfcreator={\studentname},
	pdfproducer={XeLaTeX},
	pdfkeywords={\docname}
}

% путь к каталогу с рисунками
\graphicspath{{fig/}}

% настоящее матожидание
\newcommand{\MExpect}{\mathsf{M}}

% объявили оператор!
\DeclareMathOperator{\sgn}{\mathop{sgn}}
